% Generated draft: requires mathematical and source review.
\documentclass[11pt]{article}
\usepackage[margin=25mm]{geometry}
\usepackage{fontspec}
\setmainfont{Latin Modern Roman}
\usepackage{amsmath,amssymb,mathtools,unicode-math}
\setmathfont{Latin Modern Math}
\usepackage{enumitem,needspace,xurl,hyperref}
\setlength{\parindent}{0pt}
\setlength{\parskip}{5pt}
\setlength{\emergencystretch}{3em}
\newcommand{\N}{\mathbb{N}}
\newcommand{\Z}{\mathbb{Z}}
\newcommand{\Q}{\mathbb{Q}}
\newcommand{\R}{\mathbb{R}}
\newcommand{\C}{\mathbb{C}}
\newcommand{\F}{\mathbb{F}}
\newcommand{\A}{\mathbb{A}}
\newcommand{\PP}{\mathbb{P}}
\newcommand{\E}{\mathbb{E}}
\newcommand{\eps}{\varepsilon}
\newcommand{\dd}{\,\mathrm{d}}
\begin{document}
\section*{UnsolvedMath ID: 6700080}
\textbf{Source problem code:} catalogue entry
\subsection{Definitions and mathematical statement}
In the usual closed oriented convention, an $n$-manifold $X$ is $\Q$-essential if its classifying map $c:X\to B\pi_1(X)$ satisfies $c_*[X]_{\Q}\ne0$. The source does not spell out closedness, orientation, or alternatives with local coefficients. Let $g$ be a continuous positive-definite Riemannian metric, lift it to $\widetilde X$, and use its induced length distance and Riemannian volume density. Write $\omega_n$ for the Euclidean unit-ball volume.

The proposed $C^0$-continuous Guth--Geroch lower bound is
\[
 \forall R>0\quad\exists\widetilde x_R\in\widetilde X,
 \qquad
 \operatorname{vol}_g B_{\widetilde X}(\widetilde x_R,R)
          \ge\omega_n R^n.
\]
The center may depend on $R$. The claim is not a bound for every ball or a single center working at every radius. No scalar-curvature assumption is stated, and continuous metrics are explicitly included.

The passage reports the following rough corollary of Guth's mesoscopic filling-radius theorem: for some dimension-dependent universal $\varepsilon_n>0$, the same quantifiers hold with lower bound $\varepsilon_n\omega_nR^n$. The sharp target is coefficient $1$, rather than merely a positive dimensional constant.

It then states a Burago--Ivanov asymptotic-volume rigidity assertion: if a universal cover of a $\Q$-essential $X$ has balls $\widetilde B(R_i)$ with $R_i\to\infty$ and
\[
 \operatorname{vol}(\widetilde B(R_i))\le\omega_nR_i^n,
\]
then $X$ is flat, with the radii allowed to depend on the metric. The source footnote says the cited paper concerns $X$ homeomorphic to $T^n$ and claims a reduction of the general case via polynomial-growth group classification. Those broader hypotheses and the reduction are not supplied and are retained as unverified source context, not a checked theorem covering every essential manifold.

Finally the passage claims extensions of the Burago--Ivanov argument, and of Guth's theorem, to non-Riemannian manifolds and pseudomanifolds with arbitrary metrics and suitably defined Hilbert volumes. It explicitly warns that these volumes are not a priori additive. Their definition and the generalized essentialness and volume hypotheses are absent. These extension claims are not used to strengthen the precisely displayed Riemannian conjecture.
\subsection{Short English statement}
Must every universal cover of a rationally essential manifold contain, at each radius, at least one ball with Euclidean-or-greater volume, even for a continuous metric? The source gives a rough positive-constant bound as motivation and discusses asymptotic rigidity and generalized volumes.

The supplied summary claims torus and low-dimensional progress and mentions Guth and ``Abiero/Guth,'' without exact results. It calls the general conjecture open. Those claims do not verify the all-radii sharp bound or turn the source's asymptotic discussion into a complete proof. The registry partial-progress label and 2017 passage are retained with current status marked for review.
\subsection{Sources}
\begin{itemize}
\item Local UnsolvedMath registry, permanent ID 6700080, source code catalogue entry, scalar-curvature Question [?82], proposed by Misha Gromov, 2017. Dataset curation and metadata: Ulam AI and UnsolvedMath Contributors, CC BY 4.0.
\item Misha Gromov, \emph{101 Questions, Problems and Conjectures around Scalar Curvature} (2017), printed tag [?82], occurrence 83, PDF page 82, including footnote 68, \url{https://www.ihes.fr/~gromov/wp-content/uploads/2018/08/101-problemsOct1-2017.pdf}; normalized pdf-text extraction, accessed 2026-07-29.
\item Source reference [57], conjectural bound; Guth's mesoscopic filling-radius theorem [55,57]; Burago--Ivanov asymptotic ball-volume theorem [21]; Hilbert volume [46]. Full bibliographic entries and exact theorem hypotheses are absent.
\item The research summary mentions ``Abiero/Guth'' and other recent progress without exact citations; that attribution is not corrected or verified here.
\item Source-specific provenance: the author-hosted PDF is incomplete and publicly available; current status and statement-reuse terms NEEDS\_REVIEW, separately from the dataset licence.
\end{itemize}
\end{document}
